\Needspace{9\baselineskip}
\originalchapter{含时微扰与绝热演化}
\originalsection{相互作用表象与跃迁幅}
设$H(t)=H_0+\lambda W(t)$, $H_0\ket n=E_n\ket n$. 写$\ket{\psi(t)}=\sum_n c_n(t)\e^{-\ii E_nt/\hbar}\ket n$. 代入薛定谔方程, $H_0$项与指数导数抵消, 得
\[
\ii\hbar\dot c_m(t)=\lambda\sum_nW_{mn}(t)
\e^{\ii(E_m-E_n)t/\hbar}c_n(t).
\]
若初态为$\ket i$, 零阶取$c_n=\delta_{ni}$. 代入右侧并从0到$t$积分, 得一阶跃迁幅
\[
c_f^{(1)}(t)=-\frac{\ii\lambda}{\hbar}\int_0^t
W_{fi}(s)\e^{\ii\omega_{fi}s}\dd s,
\qquad \omega_{fi}=(E_f-E_i)/\hbar.
\]
此时$f\ne i$的概率最低阶为$|c_f^{(1)}|^2$, 是$\lambda^2$阶, 不是$\lambda$阶. 只保留跃迁概率却不给初态概率作相应二阶修正, 总概率会轻微偏离1; 这是截断误差, 不是精确动力学破坏幺正性.

\originalsection{共振与黄金规则的限制}
若一段时间内 $W_{fi}(s)=w\e^{-\ii\omega s}$, 定义失谐$\delta=\omega_{fi}-\omega$, 积分给
\[
P_{i\to f}(t)\simeq\frac{4\lambda^2|w|^2}{\hbar^2}
\frac{\sin^2(\delta t/2)}{\delta^2}.
\]
$\delta=0$时取连续极限$P\simeq\lambda^2|w|^2t^2/\hbar^2$. 它在长时间会大于1, 说明一阶幅近似已超出适用范围; 不能把它当精确的永久增长规律. 对离散二态共振, 完整解产生有界的Rabi振荡.

若末态形成稠密连续谱, 记$\operatorname{sinc}z=\sin z/z$（零处连续取1）. 状态密度和矩阵元在共振窄窗内变化缓慢, 积分中的$sinc^2$核在长于相关时间且短于显著耗尽时间的窗口集中在能量守恒点. 分布意义下
$4\sin^2(\delta t/2)/\delta^2\to2\pi t\delta_{\rm D}(\delta)$.
把末态求和换为能量积分, 得Fermi黄金规则, 单位时间的跃迁概率即跃迁率为 $\Gamma=(2\pi/\hbar)\lambda^2|W_{fi}|^2\rho_E(E_i+\hbar\omega)$. 此处$W_{fi}$指驱动的对应Fourier分量$w$, 多个末态通道要分别加权求和或积分. $\rho_E$按单位能量计数, 不能与按角频率的密度混用而漏掉$\hbar$. 一个孤立的有限二态系统不满足连续谱条件, 不能直接用常数跃迁率描述全部时间.

\originalsection{绝热条件和Berry相位}
令$H(t)\ket{n(t)}=E_n(t)\ket{n(t)}$, 假定目标本征空间是一维, 与其余能级在整个路径上有非零间隙, 本征态及Hamilton算符足够平滑. 对$m\ne n$, 对本征方程求导并左乘$\bra m$, 得
$\braket m{\dot n}=\bra m\dot H\ket n/(E_n-E_m)$. 展开态后, 非对角项因快速的能隙相位振荡而抵消. 一个常用的局部小量是
\[
\frac{\hbar|\bra m\dot H\ket n|}{|E_n-E_m|^2}\ll1.
\]
它是有用的诊断, 不是不加全程光滑性、间隙和时间积累条件就成立的充分定理. 特别在避免交叉的最小能隙处, 要检查最差情况.

忽略跃迁后, 从$n$态出发的近似解为
\[
\ket{\psi(t)}\simeq
\exp\left[-\frac{\ii}{\hbar}\int_0^tE_n(s)\dd s\right]
\e^{\ii\gamma_n(t)}\ket{n(t)},
\quad \gamma_n(t)=\ii\int_0^t\braket{n(s)}{\dot n(s)}\dd s.
\]
归一化导数保证$\braket n{\dot n}$纯虚, 所以$\gamma_n$实. 第一项是动力学相位, 第二项是几何相位. 变换$\ket n\mapsto\e^{\ii\chi(t)}\ket n$会把$\gamma_n$改变$-\chi(t)+\chi(0)$, 恰好抵消本征态相位变化. 对闭合参数回路, 若采用周期规范$\ket{n(T)}=\ket{n(0)}$, Berry相位就是$\gamma_n(T)$模$2\pi$. 对任意连续规范, 规范不变的闭环量应写为 $\gamma_{\rm loop}=\gamma_n(T)+\arg\braket{n(0)}{n(T)}$模$2\pi$, 两项的规范变化恰好抵消. 这个闭环量可以通过干涉读取, 不能把任意开路径下单独的$\gamma_n$当作规范不变观测量.
